\begin{definition}[Periodic-sector projector]
    \label{def:mpdo_periodic_sector_projector}
    \lean{MPOTensor.PeriodicSectorProjector}
    \leanok
    \uses{def:mpo_tensor, def:mpo_family, def:mpdo_vertical_one_site_actions}
    Let $M$ be an MPO tensor. A \emph{periodic-sector projector} for $M$
    with exponent $p\ge 1$ is an orthogonal projector $Q$ on
    the physical space such that for every $N\ge 0$, with
    $Q_1=Q\otimes\Id^{\otimes N}$ acting on the first site,
    \begin{align}
        Q_1H^{(N+1)}
        &\ne H^{(N+1)}Q_1,
        \label{eq:mpdo_periodic_noncommute}\\
        Q_1[H^{(N+1)}]^p
        &=[H^{(N+1)}]^pQ_1.
        \label{eq:mpdo_periodic_commute_power}
    \end{align}
    These are the commutation identities that the proof of
    \cite[Proposition~4.13]{Cirac2016MPDO_arXiv}, lines 1889--1891, attaches
    to a nontrivial vertical period; the second family is the displayed
    identity at lines 1890--1891.
\end{definition}

\begin{theorem}[No periodic-sector projector for a matrix product density
    operator]
    \label{thm:mpdo_no_periodic_sector_projector}
    \lean{MPOTensor.isEmpty_periodicSectorProjector}
    \leanok
    \uses{def:mpdo, def:mpdo_periodic_sector_projector,
        thm:psd_commute_of_commute_pow}
    Let $M$ generate matrix product density operators. Then $M$ admits no
    periodic-sector projector.
\end{theorem}

\begin{proof}
    \leanok
    \uses{thm:psd_commute_of_commute_pow}
    Suppose $Q$ were such a projector. Every $H^{(N+1)}$ is positive
    semidefinite, so commutation of $Q_1$ with $[H^{(N+1)}]^p$
    implies commutation with $H^{(N+1)}$
    (Theorem~\ref{thm:psd_commute_of_commute_pow}), contradicting
    $Q_1H^{(N+1)}\ne H^{(N+1)}Q_1$. This is the contradiction concluding
    the periodic-sector step in the proof of
    \cite[Proposition~4.13]{Cirac2016MPDO_arXiv}, lines 1892--1893.
\end{proof}

\begin{definition}[Periodic vectors supply projectors]
    \label{def:mpdo_periodic_vector_yields_projector}
    \lean{MPOTensor.PeriodicVectorYieldsProjector}
    \leanok
    \uses{def:vertical_tensor_mpdo, def:no_periodic_vectors,
        def:mpdo_periodic_sector_projector}
    Say that \emph{periodic vectors of $M$ supply projectors} if for every
    nontrivial $p$-periodic vector of the vertically viewed tensor
    $\widetilde M$ --- an irreducible restriction $B$ of $\widetilde M$ to a
    minimal invariant subspace whose transfer map has a peripheral
    eigenvalue different from its spectral radius, as in
    Definition~\ref{def:no_periodic_vectors} --- there is an exponent
    $q\ge 1$ for which $M$ admits a periodic-sector projector. The proof of
    \cite[Proposition~4.13]{Cirac2016MPDO_arXiv}, lines 1889--1891, asserts
    this implication from the general theory of the peripheral spectrum of
    completely positive maps, with $q$ equal to the period $p$ of the
    vector; only the existence of some exponent enters the exclusion
    argument. The projector, its word invariance, and its single-letter
    displacement are constructed from the cyclic decomposition of the
    peripheral spectrum
    (Theorem~\ref{thm:mpdo_displaced_invariant_projector}); the passage
    from the displacement to the all-length non-commutation family is not
    needed for the periodic-sector contradiction below.
    % Gap recorded in docs/paper-gaps/cpgsv17_periodic_sector_projector.tex.
\end{definition}

\begin{theorem}[No periodic vectors in the vertical direction]
    \label{thm:mpdo_vertical_no_periodic_vectors}
    \lean{MPOTensor.hasNoPeriodicVectors_verticalTensor}
    \leanok
    \uses{def:mpdo, def:vertical_tensor_mpdo, def:no_periodic_vectors,
        def:mpdo_periodic_vector_yields_projector,
        thm:mpdo_no_periodic_sector_projector}
    Let $M$ generate matrix product density operators and suppose that
    periodic vectors of $M$ supply projectors. Then the vertically viewed
    tensor $\widetilde M$ has no nontrivial $p$-periodic vectors. This is
    the periodic-sector step in the proof of
    \cite[Proposition~4.13]{Cirac2016MPDO_arXiv}, lines 1888--1893, granted
    the projector that the source draws from the general theory of the
    peripheral spectrum.
    % Gap recorded in docs/paper-gaps/cpgsv17_periodic_sector_projector.tex.
\end{theorem}

\begin{proof}
    \leanok
    \uses{def:mpdo_periodic_vector_yields_projector,
        thm:mpdo_no_periodic_sector_projector}
    A nontrivial $p$-periodic vector of $\widetilde M$ supplies a
    periodic-sector projector for $M$, and
    Theorem~\ref{thm:mpdo_no_periodic_sector_projector} shows that no such
    projector exists.
\end{proof}

\begin{definition}[Vertically stacked layers]
    \label{def:mpo_stacked_layers}
    \lean{MPOTensor.idTensor}
    \lean{MPOTensor.layerMul}
    \lean{MPOTensor.stackedTensor}
    \leanok
    \uses{def:mpo_tensor}
    Placing one MPO layer on top of another and contracting the shared
    physical index gives the \emph{layer product} of two MPO tensors $M$ and
    $M'$ with the same physical dimension,
    $(MM')^{ij}=\sum_kM^{ik}\otimes M'^{kj}$, an MPO tensor on the tensor
    product of the virtual spaces. The \emph{$p$-fold vertical stack}
    $M^{[p]}$ of $M$ is the iterated layer product of $p$ copies of $M$; the
    empty stack $M^{[0]}$ is the trivial-bond tensor with letters
    $\delta_{ij}$, which generates the identity operator at every system
    size.
\end{definition}

\begin{theorem}[Stacked layers generate powers of the density operator]
    \label{thm:mpo_stacked_layers_power}
    \lean{MPOTensor.mpo_layerMul}
    \lean{MPOTensor.mpo_stackedTensor}
    \lean{MPOTensor.IsMPDO.stackedTensor}
    \leanok
    \uses{def:mpo_stacked_layers, def:mpo_family, def:mpdo}
    For every system size $N$, the operator generated by a layer
    product is the product of the density operators generated by the layers,
    $H^{(N)}(MM')=H^{(N)}(M)\,H^{(N)}(M')$.
    Consequently the $p$-fold vertical stack generates
    $H^{(N)}(M^{[p]})=[H^{(N)}(M)]^p$ for all $N$ and $p$; this is the
    reading of the powers in the displayed identity of
    \cite[Proposition~4.13, lines~1890--1891]{Cirac2016MPDO_arXiv} as
    density operators of stacked layers. In particular, the vertical stack
    of a tensor generating matrix product density operators again generates
    matrix product density operators on every nonempty chain.
\end{theorem}

\begin{proof}
    \leanok
    \uses{def:mpo_stacked_layers, def:mpo_family}
    The word product of layer-product letters along a ket--bra configuration
    expands, by induction on the length and the mixed-product property of
    the Kronecker product, as a sum over configurations $\kappa$ of the
    contracted physical index of Kronecker products of the layers' word
    products. Taking traces, each summand factors into the product of the
    two traces, so the entry of $H^{(N)}(MM')$ at $(\sigma,\tau)$ equals
    $\sum_\kappa H^{(N)}(M)_{\sigma\kappa}H^{(N)}(M')_{\kappa\tau}$, the
    entry of the matrix product. The power identity follows by induction on
    $p$, and positive semidefiniteness of $[H^{(N)}]^p$ follows from that of
    $H^{(N)}$.
\end{proof}

\begin{theorem}[Vertical letters of the stacked tensor]
    \label{thm:mpo_stacked_vertical_letters}
    \lean{MPOTensor.verticalTensor_layerMul}
    \lean{MPOTensor.exists_word_verticalTensor_stackedTensor}
    \lean{MPOTensor.stackedTensor_ketLeftMul_invariant}
    \leanok
    \uses{def:mpo_stacked_layers, def:vertical_tensor_mpdo,
        def:mpdo_vertical_one_site_actions}
    The vertically viewed letters of a layer product are the products of the
    layers' vertically viewed letters,
    $(MM')_{(a,a'),(b,b')}=M_{ab}M'_{a'b'}$. Hence every vertically viewed
    letter of the $p$-fold stack $M^{[p]}$ is the evaluation of a length-$p$
    word of $\widetilde M$. In particular, if a matrix $Q$ satisfies the
    one-sided invariance $QA_w=QA_wQ$ for every product $A_w$ of $p$ letters
    of $\widetilde M$, then $QM^{[p]}=QM^{[p]}Q$ in the sense of
    Definition~\ref{def:mpdo_vertical_one_site_actions}.
\end{theorem}

\begin{proof}
    \leanok
    \uses{def:mpo_stacked_layers, def:vertical_tensor_mpdo,
        def:mpdo_vertical_one_site_actions}
    The letter identity is a direct computation from the layer product:
    contracting the shared physical index composes the physical-space
    operators. Induction on $p$ writes each vertical letter of $M^{[p]}$ as
    a length-$p$ word. Thus, for every bond index $v$, there is a word $w$
    of length $p$ such that $\widetilde{M^{[p]}}_v=A_w$, and
    $QA_w=QA_wQ$ implies
    $Q\widetilde{M^{[p]}}_v=Q\widetilde{M^{[p]}}_vQ$.
    Applying this identity to the two one-site actions gives
    $QM^{[p]}=QM^{[p]}Q$.
\end{proof}

\begin{definition}[Periodic-sector projector from cyclic data]
    \label{def:mpdo_periodic_sector_projector_of_cyclic}
    \lean{MPOTensor.periodicSectorProjectorOfCyclicData}
    \leanok
    \uses{def:mpdo, def:mpdo_periodic_sector_projector,
        def:mpdo_vertical_one_site_actions, def:mpo_stacked_layers}
    Let $M$ generate matrix product density operators, let $p\ge 1$, and let
    $Q$ be an orthogonal projector on the physical space satisfying
    $QA_w=QA_wQ$ for every product $A_w$ of $p$ letters of $\widetilde M$
    and $Q_1H^{(N)}\ne H^{(N)}Q_1$ for every length. Then $Q$ is a
    periodic-sector projector for $M$ with exponent $p$.
\end{definition}

\begin{proof}
    \leanok
    \uses{thm:mpo_stacked_vertical_letters, thm:mpo_stacked_layers_power,
        thm:mpdo_first_site_invariant_comm}
    By Theorem~\ref{thm:mpo_stacked_vertical_letters} the word invariance
    gives $QM^{[p]}=QM^{[p]}Q$ on the stacked tensor, which generates matrix
    product density operators by
    Theorem~\ref{thm:mpo_stacked_layers_power}.
    Theorem~\ref{thm:mpdo_first_site_invariant_comm} applied to $M^{[p]}$
    yields $Q_1H^{(N)}(M^{[p]})=H^{(N)}(M^{[p]})Q_1$ for every length, and
    $H^{(N)}(M^{[p]})=[H^{(N)}]^p$ by the stacked-layers identity. Together
    with the assumed non-commutation family this is exactly the
    periodic-sector projector of
    Definition~\ref{def:mpdo_periodic_sector_projector}.
\end{proof}

\begin{definition}[Periodic vectors supply cyclic projectors]
    \label{def:mpdo_periodic_vector_yields_cyclic_projector}
    \lean{MPOTensor.PeriodicVectorYieldsCyclicProjector}
    \leanok
    \uses{def:vertical_tensor_mpdo, def:no_periodic_vectors,
        def:mpdo_vertical_one_site_actions}
    Say that \emph{periodic vectors of $M$ supply cyclic projectors} if for
    every nontrivial $p$-periodic vector of the vertically viewed tensor
    $\widetilde M$ --- presented as in
    Definition~\ref{def:no_periodic_vectors} --- there are an exponent
    $q\ge 1$ and an orthogonal projector $Q$ on the physical space
    satisfying the one-sided invariance $QA_w=QA_wQ$ for every product
    $A_w$ of $q$ letters of $\widetilde M$, together with
    $Q_1H^{(N)}\ne H^{(N)}Q_1$ for every length. This is the shape of the
    cyclic decomposition of the peripheral spectrum
    \cite[Theorem~6.6]{Wolf2012Quantum} applied to the vertical transfer
    map: each cyclic projection of a period-$p$ sector returns to itself
    after $p$ applications of the map, while its nontrivial cyclic
    displacement under one application prevents commutation with
    $H^{(N)}$. Compared with
    Definition~\ref{def:mpdo_periodic_vector_yields_projector}, the
    commutation of $Q$ with $[H^{(N)}]^p$ is no longer assumed; it is
    derived through the stacked-layers identity.
\end{definition}

\begin{theorem}[Reduction of the periodic-sector step to the cyclic
    decomposition]
    \label{thm:mpdo_cyclic_projector_reduction}
    \lean{MPOTensor.periodicVectorYieldsProjector_of_cyclic}
    \lean{MPOTensor.hasNoPeriodicVectors_verticalTensor_of_cyclicProjector}
    \leanok
    \uses{def:mpdo, def:mpdo_periodic_vector_yields_cyclic_projector,
        def:mpdo_periodic_vector_yields_projector,
        def:vertical_tensor_mpdo, def:no_periodic_vectors}
    Let $M$ generate matrix product density operators. If periodic vectors
    of $M$ supply cyclic projectors, then they supply periodic-sector
    projectors, and the vertically viewed tensor $\widetilde M$ has no
    nontrivial $p$-periodic vectors. This reduces the hypothesis of the
    periodic-sector step in the proof of
    \cite[Proposition~4.13, lines~1888--1893]{Cirac2016MPDO_arXiv} to the
    cyclic decomposition of the peripheral spectrum of the vertical
    transfer map.
\end{theorem}

\begin{proof}
    \leanok
    \uses{def:mpdo_periodic_sector_projector_of_cyclic,
        thm:mpdo_vertical_no_periodic_vectors}
    A nontrivial $p$-periodic vector supplies an invariant non-commuting
    projector, which is a periodic-sector projector by
    Definition~\ref{def:mpdo_periodic_sector_projector_of_cyclic}. The
    exclusion of periodic vectors is then
    Theorem~\ref{thm:mpdo_vertical_no_periodic_vectors}.
\end{proof}

\begin{theorem}[The displaced invariant projector of a nontrivial periodic
    vector]
    \label{thm:mpdo_displaced_invariant_projector}
    \lean{MPOTensor.exists_displaced_invariant_projector_of_periodic_vector}
    \leanok
    \uses{def:vertical_tensor_mpdo, def:no_periodic_vectors,
        def:mpdo_vertical_one_site_actions, def:support_proj}
    Let $M$ be an MPO tensor and let $\widetilde M$ carry a nontrivial
    $p$-periodic vector, presented as in
    Definition~\ref{def:no_periodic_vectors}: an irreducible restriction
    $B$ of $\widetilde M$ to a minimal invariant subspace, reached through
    an isometry $V$, with a positive definite transfer eigenvector $\rho$
    of eigenvalue $r>0$ and a transfer eigenvalue $\mu\ne r$ of modulus
    $r$. Then there are an exponent $q\ge 1$ and an orthogonal projector
    $Q$ on the physical space such that $QA_w=QA_wQ$ for every product
    $A_w$ of $q$ letters of $\widetilde M$, while
    $Q\widetilde M\ne Q\widetilde MQ$ in the sense of
    Definition~\ref{def:mpdo_vertical_one_site_actions}. This realizes
    the projector that the proof of
    \cite[Proposition~4.13, lines~1889--1891]{Cirac2016MPDO_arXiv} draws
    from the general theory of the peripheral spectrum.
\end{theorem}

\begin{proof}
    \leanok
    \uses{thm:spectral_unital_gauge_schwarz_setup,
        lem:eigenvalue_transport_unital_gauge,
        thm:peripheral_cyclic_structure,
        thm:cyclic_decomposition_irreducible_schwarz,
        lem:kraus_cyclic_shift, lem:support_proj_range_factorization}
    The rescaled unital gauge $K$ of the corner satisfies the hypotheses of
    the cyclic decomposition of the peripheral spectrum
    (Theorem~\ref{thm:spectral_unital_gauge_schwarz_setup}), and the
    eigenvalue $\mu/r\ne 1$ of modulus one
    (Lemma~\ref{lem:eigenvalue_transport_unital_gauge}) forces the
    peripheral cyclic group to have order $m\ge 2$.
    Theorem~\ref{thm:cyclic_decomposition_irreducible_schwarz} produces
    orthogonal projections $P_0,\ldots,P_{m-1}$ summing to the identity
    with $\E_K(P_{k+1})=P_k$, and the letter-level shift
    $K_vP_{k+1}=P_kK_v$ holds
    (Lemma~\ref{lem:kraus_cyclic_shift}). Set $q=m$ and let $Q$ be the
    complement of the support projection $\pi$ of
    $YY^\dagger$ for $Y=V\rho^{1/2}P_0$. A word $A_w$ of $m$ letters of
    $\widetilde M$ satisfies $A_wY=YG_w$, where $G_w$ is a scalar multiple
    of the corresponding word in the letters of $K$, which commutes with
    $P_0$ by the full-period shift; so the range of $Y$ is invariant and
    Lemma~\ref{lem:support_proj_range_factorization} gives
    $QA_w=QA_wQ$. If a single letter also satisfied the invariance, then
    $\pi$ would fix $A_vY$, whose columns lie in the sector
    $V\rho^{1/2}P_{m-1}$ by the cyclic shift; cancelling the injective
    factor $V\rho^{1/2}$ and using orthogonality of $P_{m-1}$ and $P_0$
    would force $P_{m-1}K_v=0$ for every letter, hence
    $P_{m-1}=\E_K(P_0)=0$, contradicting the nonvanishing of the cyclic
    projections.
\end{proof}

\begin{theorem}[A displaced idempotent fails to commute at some length]
    \label{thm:cpsv_literal_displaced_noncommutation}
    \lean{MPSTensor.IsCPSVCanonicalForm.exists_not_commute_of_displaced}
    \leanok
    \uses{def:cpsv_canonical_form, def:mpo_family,
        def:mpdo_vertical_one_site_actions}
    Let $M$ be an MPO tensor whose doubled-index MPS tensor is in literal
    CPSV canonical form. If $Q^2=Q$ and
    $Q\widetilde M\ne Q\widetilde M Q$, then there is an integer $N\geq0$
    such that, for $Q_1=Q\otimes\Id^{\otimes N}$,
    \begin{align}
        Q_1H^{(N+1)}&\ne H^{(N+1)}Q_1.
        \label{eq:cpsv_literal_displaced_noncommutation}
    \end{align}
    Thus a displaced idempotent has a noncommuting chain length. The
    proof of \cite[Proposition~4.13, lines~1888--1893]{Cirac2016MPDO_arXiv}
    only requires this existential conclusion, not noncommutation at every
    length.
\end{theorem}

\begin{proof}
    \leanok
    \uses{thm:representative_one_sub_mp_zero,
        lem:mpdo_first_site_insertions_vertical_mul}
    Suppose instead that $Q_1H^{(N+1)}=H^{(N+1)}Q_1$ for every $N\geq0$.
    Idempotence gives
    \begin{align}
        Q_1H^{(N+1)}
        &=Q_1^2H^{(N+1)}
         =Q_1H^{(N+1)}Q_1.
        \notag
    \end{align}
    Taking matrix coefficients yields equality of the one-sided first-site
    contractions. The literal form of Lemma~L in
    Theorem~\ref{thm:representative_one_sub_mp_zero} identifies the two
    inserted tensors on the original bond space. By
    Lemma~\ref{lem:mpdo_first_site_insertions_vertical_mul}, this equality is
    $Q\widetilde M=Q\widetilde M Q$, contradicting the displacement.
\end{proof}

\begin{theorem}[Literal canonical form excludes vertical periodic sectors]
    \label{thm:mpdo_vertical_no_periodic_vectors_literal_cpsv}
    \lean{MPOTensor.hasNoPeriodicVectors_verticalTensor_of_cpsvCanonicalForm}
    \leanok
    \uses{def:mpdo, def:cpsv_canonical_form, def:vertical_tensor_mpdo,
        def:no_periodic_vectors}
    Let $M$ generate matrix product density operators, and suppose that its
    doubled-index MPS tensor is in literal CPSV canonical form. Then the
    vertically viewed tensor $\widetilde M$ has no nontrivial periodic
    vectors.
\end{theorem}

\begin{proof}
    \leanok
    \uses{thm:mpdo_displaced_invariant_projector,
        thm:cpsv_literal_displaced_noncommutation,
        thm:mpo_stacked_vertical_letters,
        thm:mpo_stacked_layers_power,
        thm:mpdo_first_site_invariant_comm,
        thm:psd_commute_of_commute_pow}
    A nontrivial periodic vector supplies an integer $p\geq1$ and an
    orthogonal projector $Q$ such that every word $A_w$ of length $p$ in the
    letters of $\widetilde M$ satisfies $QA_w=QA_wQ$, while
    $Q\widetilde M\ne Q\widetilde M Q$
    (Theorem~\ref{thm:mpdo_displaced_invariant_projector}). By
    Theorem~\ref{thm:cpsv_literal_displaced_noncommutation}, there is a length
    $N+1$ for which
    \begin{align}
        Q_1H^{(N+1)}&\ne H^{(N+1)}Q_1.
        \notag
    \end{align}
    The word invariance and the vertical-stack identity give
    \begin{align}
        Q_1[H^{(N+1)}]^p&=[H^{(N+1)}]^pQ_1.
        \notag
    \end{align}
    Since $H^{(N+1)}$ is positive semidefinite,
    Theorem~\ref{thm:psd_commute_of_commute_pow} implies
    $Q_1H^{(N+1)}=H^{(N+1)}Q_1$, a contradiction. This is the periodic
    exclusion in
    \cite[Proposition~4.13, lines~1888--1893]{Cirac2016MPDO_arXiv}, with the
    single noncommuting length supplied by Lemma~L.
\end{proof}

\begin{theorem}[BNT-refined horizontal form excludes vertical periodic sectors]
    \label{thm:mpdo_vertical_no_periodic_vectors_horizontal_cf}
    \lean{MPOTensor.IsHorizontalCF.exists_not_commute_of_displaced}
    \lean{MPOTensor.hasNoPeriodicVectors_verticalTensor_of_horizontalCF}
    \leanok
    \uses{def:mpdo, def:horizontal_cf_mpdo, def:vertical_tensor_mpdo,
        def:no_periodic_vectors, thm:representative_one_sub_mp_zero,
        lem:mpdo_first_site_insertions_vertical_mul,
        thm:mpdo_displaced_invariant_projector,
        thm:mpo_stacked_layers_power, thm:mpdo_first_site_invariant_comm,
        thm:psd_commute_of_commute_pow}
    Let $M$ be in normalized BNT-refined horizontal form. If an idempotent
    $Q$ satisfies
    $Q\widetilde M\ne Q\widetilde MQ$, then for some $N\ge 0$,
    $Q_1H^{(N+1)}\ne H^{(N+1)}Q_1$.
    Consequently, if $M$ generates matrix product density operators, then
    $\widetilde M$ has no nontrivial periodic vectors.

    \textbf{Scope restriction (BNT-refined horizontal form):}
    the normalized BNT-refined horizontal hypothesis is stronger than the
    literal CPSV canonical-form hypothesis.
    % The literal canonical-form conclusion is
    % Theorem~\ref{thm:mpdo_vertical_no_periodic_vectors_literal_cpsv}.
\end{theorem}

\begin{proof}
    \leanok
    \uses{thm:representative_one_sub_mp_zero,
        lem:mpdo_first_site_insertions_vertical_mul,
        thm:mpdo_displaced_invariant_projector,
        thm:mpo_stacked_layers_power, thm:mpdo_first_site_invariant_comm,
        thm:psd_commute_of_commute_pow}
    Suppose first that $Q_1$ commutes with $H^{(N+1)}$ for every $N$.
    Since $Q^2=Q$, this gives
    $Q_1H^{(N+1)}=Q_1H^{(N+1)}Q_1$ at every positive length.
    Theorem~\ref{thm:representative_one_sub_mp_zero} and
    Lemma~\ref{lem:mpdo_first_site_insertions_vertical_mul} then imply
    $Q\widetilde M=Q\widetilde MQ$, a contradiction. Thus a displaced
    idempotent has a noncommuting length $N+1$.

    A nontrivial periodic vector supplies the orthogonal projector $Q$ of
    Theorem~\ref{thm:mpdo_displaced_invariant_projector}, invariant under
    every word of some full period $p\ge1$. At the noncommuting length just
    obtained, the stacked-layer identity and the invariant-projection
    calculation give
    $Q_1[H^{(N+1)}]^p=[H^{(N+1)}]^pQ_1$.
    Positive semidefiniteness and
    Theorem~\ref{thm:psd_commute_of_commute_pow} imply that $Q_1$ commutes
    with $H^{(N+1)}$, which is the required contradiction. This uses the
    pointwise contradiction in
    \cite[Proposition~4.13, lines~1888--1893]{Cirac2016MPDO_arXiv}; the
    stronger all-length inequality printed there is unnecessary.
\end{proof}

\begin{definition}[Displaced projectors never commute]
    \label{def:mpdo_noninvariant_projector_noncommuting}
    \lean{MPOTensor.NoninvariantProjectorNoncommuting}
    \leanok
    \uses{def:mpo_tensor, def:mpo_family, def:mpdo_vertical_one_site_actions}
    Say that \emph{displaced projectors of $M$ never commute} if every
    orthogonal projector $Q$ on the physical space with
    $Q\widetilde M\ne Q\widetilde MQ$ satisfies
    $Q_1H^{(N)}\ne H^{(N)}Q_1$ at every length. This is the
    canonical-form step of the periodic-sector argument: the proof of
    \cite[Proposition~4.13]{Cirac2016MPDO_arXiv} assumes the tensor is in
    canonical form in the horizontal direction, and the argument of
    lines~1874--1887 turns all-length commutation of $Q_1$ with $H^{(N)}$
    into the letter-level invariance $Q\widetilde M=Q\widetilde MQ$. This
    argument uses the normalized BNT-refined horizontal hypothesis. The
    all-length condition is stronger than is needed for either
    Theorem~\ref{thm:mpdo_vertical_no_periodic_vectors_literal_cpsv} or
    Theorem~\ref{thm:mpdo_vertical_no_periodic_vectors_horizontal_cf}; each
    uses a single noncommuting length supplied by the contrapositive of
    Lemma~L.
\end{definition}

\begin{theorem}[Reduction of the periodic-sector step to the
    non-commutation family]
    \label{thm:mpdo_cyclic_projector_of_noncommutation}
    \lean{MPOTensor.periodicVectorYieldsCyclicProjector_of_noncommutation}
    \lean{MPOTensor.hasNoPeriodicVectors_verticalTensor_of_noncommutation}
    \leanok
    \uses{def:mpdo, def:mpdo_noninvariant_projector_noncommuting,
        def:mpdo_periodic_vector_yields_cyclic_projector,
        def:vertical_tensor_mpdo, def:no_periodic_vectors}
    Let $M$ generate matrix product density operators. If displaced
    projectors of $M$ never commute, then periodic vectors of $M$ supply
    cyclic projectors, and the vertically viewed tensor $\widetilde M$ has
    no nontrivial $p$-periodic vectors. This narrows the hypothesis of the
    periodic-sector step in the proof of
    \cite[Proposition~4.13, lines~1888--1893]{Cirac2016MPDO_arXiv} to the
    canonical-form implication of lines~1874--1887: every other ingredient
    of the step is proved.
\end{theorem}

\begin{proof}
    \leanok
    \uses{thm:mpdo_displaced_invariant_projector,
        thm:mpdo_cyclic_projector_reduction}
    A nontrivial $p$-periodic vector supplies the invariant, displaced
    projector of Theorem~\ref{thm:mpdo_displaced_invariant_projector}; the
    non-commutation hypothesis upgrades the displacement to the all-length
    non-commutation family, so periodic vectors supply cyclic projectors
    and Theorem~\ref{thm:mpdo_cyclic_projector_reduction} applies.
\end{proof}

\begin{theorem}[Commutation at chain length two forces letter invariance, for
    an injective tensor]
    \label{thm:mpdo_commute_of_isInjective}
    \lean{MPOTensor.ketLeftMul_eq_braRightMul_of_commute_of_isInjective}
    \leanok
    \uses{def:injective, def:mpdo_vertical_one_site_actions, def:mpo_family}
    Let $M$ be an MPO tensor whose doubled-index tensor is injective, and
    let $Q$ be an idempotent on the physical space (Hermiticity is not
    needed). If $Q_1H^{(2)}=H^{(2)}Q_1$, then
    $Q\widetilde M=Q\widetilde MQ$.
\end{theorem}

\begin{proof}
    \leanok
    Idempotence of $Q$ upgrades the commutation to the one-sided reduction
    $Q_1H^{(2)}=Q_1H^{(2)}Q_1$. Unfolding this identity against every
    trailing letter of $M$ gives, for every pair of physical indices $a,b$,
    a trace identity pairing $(Q\widetilde M)_{ab}$ against
    $(Q\widetilde MQ)_{ab}$ over every letter of the doubled-index tensor.
    Injectivity's spanning property and nondegeneracy of the trace pairing
    force the two matrices to agree.
\end{proof}

\begin{theorem}[No periodic vectors for an injective matrix product density
    operator]
    \label{thm:mpdo_vertical_no_periodic_vectors_injective}
    \lean{MPOTensor.hasNoPeriodicVectors_verticalTensor_of_isInjective}
    \leanok
    \uses{def:mpdo, def:injective, def:vertical_tensor_mpdo,
        def:no_periodic_vectors, thm:mpdo_displaced_invariant_projector,
        thm:mpdo_commute_of_isInjective}
    Let $M$ generate matrix product density operators and let $M$'s
    doubled-index tensor be injective. Then the vertically viewed tensor
    $\widetilde M$ has no nontrivial periodic vectors, unconditionally.
\end{theorem}

\begin{proof}
    \leanok
    A nontrivial periodic vector supplies the invariant, displaced projector
    $Q$ of Theorem~\ref{thm:mpdo_displaced_invariant_projector}. Its word
    invariance transfers to the stacked tensor, giving commutation with
    $[H^{(2)}]^p$ unconditionally through the invariant-projection step
    applied to the stack; positive semidefiniteness of the density operators
    removes the power, and Theorem~\ref{thm:mpdo_commute_of_isInjective}
    turns the resulting commutation with $H^{(2)}$ into the letter-level
    invariance that the displacement forbids.

    Theorem~\ref{thm:mpdo_vertical_no_periodic_vectors_horizontal_cf}
    gives the corresponding result for normalized BNT-refined horizontal form.
\end{proof}
